\documentclass[
    language=english,
    doctype=article,
    institution=none,
]{../../omnilatex}

\RequirePackage{config/document-settings}
\addbibresource{../../bib/bibliography.bib}

\begin{document}

\maketitle

\section{Product Overview}

The Acme Platform v3.1 introduces the Advanced Workflow Automation module,
a visual, low-code environment that enables enterprise customers to design,
deploy, and manage complex operational workflows without writing custom code.
This module builds on the real-time analytics capabilities introduced in v3.0
and integrates with the existing API layer, identity management, and
monitoring infrastructure.

The primary target users are operations teams, DevOps engineers, and
line-of-business managers who need to automate repetitive processes, enforce
compliance rules, and orchestrate cross-functional workflows.

\section{Features}

\begin{table}[h]
    \centering
    \begin{tabular}{llp{5.5cm}}
        \toprule
        \textbf{Feature} & \textbf{Priority} & \textbf{Description} \\
        \midrule
        Visual Workflow Builder    & P0 & Drag-and-drop canvas for designing
            multi-step workflows with conditional branching, loops, and
            error handling. \\
        Pre-built Connectors      & P0 & Native integrations with 50+
            SaaS and infrastructure services (AWS, Azure, GCP, Slack,
            Jira, Salesforce, etc.). \\
        Conditional Logic Engine   & P0 & Rule-based branching using
            JSONPath expressions, time-based triggers, and webhook
            events. \\
        Audit Trail                & P1 & Immutable log of all workflow
            executions, state changes, and user actions for compliance
            and debugging. \\
        Role-Based Access Control  & P1 & Granular permissions for
            workflow creation, execution, and modification at the
            folder and workflow level. \\
        Version Control            & P1 & Full version history with
            diff view, rollback, and branch/merge capabilities for
            workflow definitions. \\
        Template Library           & P2 & Curated collection of 25+
            industry-specific workflow templates covering ITSM, HR
            onboarding, and data pipeline orchestration. \\
        Advanced Analytics         & P2 & Real-time dashboards for
            workflow performance, error rates, and resource
            utilization with exportable reports. \\
        \bottomrule
    \end{tabular}
    \caption{Feature matrix for Acme Platform v3.1.}
\end{table}

\section{Requirements}

\subsection{Functional Requirements}

\begin{enumerate}
    \item The workflow builder shall support workflows with up to 200 steps
          and 50 conditional branches without performance degradation.
    \item All workflow executions shall complete within the configured
          timeout period, with a default timeout of 30 minutes.
    \item The system shall support at least 1,000 concurrent workflow
          executions per tenant.
    \item Workflow definitions shall be exportable as JSON for portability
          and version control.
    \item The audit trail shall retain records for a minimum of 7 years
          in compliance with enterprise retention policies.
\end{enumerate}

\subsection{Non-Functional Requirements}

\begin{itemize}
    \item \textbf{Performance:} Workflow execution latency shall not exceed
          500ms for the first step and 200ms for subsequent steps under
          normal load.
    \item \textbf{Availability:} The workflow engine shall maintain 99.99\%
          uptime, measured monthly.
    \item \textbf{Scalability:} The system shall support horizontal scaling
          to handle a 10x increase in workflow volume without manual
          intervention.
    \item \textbf{Security:} All workflow data in transit and at rest shall
          be encrypted using AES-256. OAuth 2.0 and SAML 2.0 shall be
          supported for connector authentication.
    \item \textbf{Compliance:} The module shall comply with SOC 2 Type II,
          GDPR, and CCPA requirements.
\end{itemize}

\section{Architecture}

The workflow automation module is composed of three primary subsystems:

\begin{enumerate}
    \item \textbf{Workflow Orchestrator:} A state machine engine built on
          an event-sourced architecture. It processes workflow definitions,
          manages execution state, and coordinates connector invocations.
          Deployed as a horizontally scalable microservice on Kubernetes.
    \item \textbf{Connector Runtime:} An isolated execution environment for
          third-party integrations. Each connector runs in a sandboxed
          container with resource limits to prevent noisy-neighbor issues.
          Connectors communicate with the orchestrator via gRPC.
    \item \textbf{Persistence Layer:} Workflow definitions and execution
          history are stored in PostgreSQL with read replicas for analytics
          queries. Audit logs are streamed to Apache Kafka for long-term
          retention in object storage.
\end{enumerate}

\begin{figure}[h]
    \centering
    \f\boxed{\par\boxed{0.85\textwidth}{\centering\vspace{2cm}
    \textbf{Architecture Diagram Placeholder}\\[0.5em]
    Workflow Orchestrator $\leftrightarrow$ Connector Runtime $\leftrightarrow$ Persistence Layer
    \vspace{2cm}}}
    \caption{High-level architecture of the workflow automation module.}
\end{figure}

\section{Timeline}

\begin{table}[h]
    \centering
    \begin{tabular}{llr}
        \toprule
        \textbf{Milestone} & \textbf{Deliverable} & \textbf{Target Date} \\
        \midrule
        M1: Design Complete     & Architecture review, API specs        & July 15, 2026  \\
        M2: Core Engine         & Workflow orchestrator, basic UI       & September 1, 2026 \\
        M3: Connectors          & 25 pre-built connectors              & October 15, 2026 \\
        M4: Beta Release        & Feature-complete beta for 10 tenants  & November 30, 2026 \\
        M5: GA Release          & Production launch with full support   & January 15, 2027 \\
        \bottomrule
    \end{tabular}
    \caption{Development milestones for Acme Platform v3.1.}
\end{table}

\section{Success Metrics}

Following the general availability release, the following metrics will be
tracked to evaluate adoption and impact:

\begin{itemize}
    \item \textbf{Adoption:} 40\% of enterprise tenants activate the
          workflow module within 90 days of GA.
    \item \textbf{Engagement:} Average of 15 active workflows per tenant
          within 6 months.
    \item \textbf{Efficiency:} Customers report a 25\% reduction in manual
          operational tasks within the first quarter of use.
    \item \textbf{Reliability:} Workflow execution success rate above 99.5\%.
    \item \textbf{Revenue:} \$8M in incremental ARR attributable to the
          workflow module in FY2027.
\end{itemize}

\printbibliography[heading=none]

\end{document}
